\subsection{Cross-Platform Considerations}
\label{sec:cross-platform}

\FractalShark{} is developed primarily as a Windows application built with
Visual~Studio~2026, MSBuild, and the NVIDIA CUDA toolkit.  Linux is supported
as a secondary native target built with CMake and Clang, with release binaries
targeted at Ubuntu.  The Linux build includes the CPU-side numeric core, CUDA
rendering kernels, GPU high-precision arithmetic, command-line renderer, unit
tests, and the native \code{FractalSharkGuiLinux} graphical application built
on Xlib and Dear~ImGui.  This subsection documents the parallel build
infrastructure, the platform-abstraction layer, and the cross-platform
serialization contract for arbitrary-precision integers.  The project goal is
full Linux/Win32 feature parity, while Windows remains the historical primary
development environment.

\subsubsection{Dual build system}
\label{subsec:dual-build}

The Windows build is authoritative.  All project structure, source membership,
and compiler settings are tracked first in MSBuild project files
(\code{.vcxproj}) and a Visual~Studio solution (\code{.sln}).  The Linux build
mirrors the native Linux target set using \code{CMakeLists.txt} files placed
alongside the corresponding \code{.vcxproj}.  Projects that remain Windows-only
have no CMake file at all.  When a shared or native-Linux project gains or
loses sources, both build systems must be updated together.

The Linux toolchain is Clang only; GCC is not a supported host compiler.  CUDA
device compilation uses the standard NVIDIA \code{nvcc} driver, with a
host-compiler split where required: \code{clang++} drives ordinary C++
compilation, while \code{nvcc} dispatches to a system \code{g++} purely as
its host compiler for device translation units.  This arrangement isolates
toolchain incompatibilities to the device-compilation step and keeps the host
C++ surface uniformly Clang-compiled.

Project-by-project coverage on the Linux side includes the shared platform
library, core fractal library, CUDA rendering library, GPU high-precision
arithmetic library, explicit-instantiation generator, GPU test harness,
portable CPU unit tests, command-line renderer, and native Linux GUI.  Legacy
Windows utilities and the Win32 GUI project remain MSBuild-only; the Linux GUI
entry point is \code{FractalSharkGuiLinux}.

Continuous integration runs the Windows MSBuild and the Linux CMake build as
separate jobs in the same matrix workflow, alongside a third job that compiles
this document.  Both builds must succeed before any change merges, ensuring the
two build systems remain synchronized as the codebase evolves.

\subsubsection{Platform abstraction layer}
\label{subsec:platform-abstraction}

The \code{FractalSharkPlatform} library provides a thin abstraction over the
small number of operating-system primitives that genuinely differ between
Windows and Linux.  Its surface comprises an \code{Environment} facade with
operations for system-heap allocation, aligned allocation, and registration of
heap-cleanup hooks invoked at process shutdown.  Each operation has separate
Windows and Linux translation units that supply the platform-specific
implementation while presenting an identical C++ interface to callers.

The abstraction is intentionally minimal.  Rather than scatter
\code{\#ifdef~\_WIN32} guards through the codebase, divergent behavior is
concentrated in a small number of files whose entire contents are
platform-specific.  Code outside the abstraction layer therefore compiles
unmodified on both platforms.  The audit policy tolerates legacy
platform-conditional sites that predate the abstraction, but new platform
divergence is expected to flow through \code{FractalSharkPlatform} unless the
divergence is truly local in scope.

\subsubsection{Arbitrary-precision serialization}
\label{subsec:mpir-portability}

\FractalShark{} relies on MPIR for arbitrary-precision integer arithmetic on
Windows and on GMP for the same role on Linux.  The two libraries are
source-compatible at the API level but differ in some internal representation
details, which matters because reference-orbit files and saved views contain
serialized arbitrary-precision integers that must round-trip across platforms.

To guarantee byte-identical files, \FractalShark{} defines a single canonical
wire format that matches MPIR's native \code{\_\_mpz\_struct} layout.  On
Windows this format is produced by direct field access into the underlying
MPIR object.  On Linux the same byte sequence is reconstructed atop GMP using
\code{mpz\_export} and \code{mpz\_import}, which expose limb data through a
portable interface.  The two implementations are validated against a checked-in
golden file (\code{golden\_mpir.bin}) by the \code{TestMpirSerialization}
unit test, which runs on both platforms in CI.  Any divergence in the wire
format fails the test on whichever platform regressed, preventing silent
incompatibility between Windows-written and Linux-written orbit files.

